\section{Principle IV: Authority}

\begin{quotex}
From Chapter II, ``Principles", L'Autorité, of \textit{Mes Idées Politiques}, by Charles Maurras.
\end{quotex}

\paragraph{Nature of Authority}
The idea that authority can be constructed from below would not have come into the head of our grandparents who were wise.

It is not made, in truth, from below nor from above.

Authority is \emph{born}. In individuals, families, peoples, it is a gift which the will of men has very little to see.

The most vulgar observation is totally in accord here with the Catholic text \emph{omnis potestas a Deo}\footnote{All power comes from God.}. In one of the oldest ``Letters to Françoise", Mr. Marcel Prevost reminds his niece how, after having passed her baccalaureate, she knew she was obliged to show the secretary of the Faculty her notes, regardless of the rules of the University and the laws of the State. ``You have a great gift, Françoise, it is authority," observed Uncle Marcel sententiously. Did he believe so well what he observed. Did he follow through what he observed?

Authority, seized thus at its birth, is something simple and pure. Certain human types possess it, the others are devoid of it.

In leaving aside those who know only how to submit, the man of liberty, recognizable in the pride of a heart that nothing crushes, differs from the one who is characterized by dignity and inspires above all respect: the man of authority differs from the other two. His liberty imposes itself naturally on the liberty of others, his dignity is radiant, it leads and carries. It is neither respect nor admiration—inert feelings—it is an enthusiastic submissiveness which responds to him.

Far from being irrational, instinctive wishes move more quickly than conscious reason, and perceptive logic is not more lacking than the passions of a great love. The author of La Vita Nuova tells us that at his first sight of Beatrice his heart begins to beat in him impetuously, which \textbf{Dante} develops and explains in these terms:

\begin{quotex}
The spirit of life that reside in the most secret vault of the heart begins to tremble with so much force that the movement is made to be felt again in my smallest veins and, trembling, it says these words: `\emph{Ecce deus fortiori me, qui veniens dominabitur me}.' Here is that God stronger than I, he is going to dominate me. Then, the animal spirit, which is hidden in its high vault where all sensitive spirits go to carry their perception, began to be greatly surprised and, addressing itself particularly to the spirits of the sight, says these words: `\emph{Apparuit jam beatitude nostra}.' Our bliss has appeared. 

\end{quotex}
It is necessary to reread all that penetrating and poetic analysis which is of an age where the laziest sophisms from Germany and the Jews [Freud] had not imposed on the European West a ridiculous philosophy of the unconscious. What was unconscious, one brought it to consciousness. What eludes the first grasp of reason, a more subtle reason pulls out in the night.

This explanation of the strong premonitions of a loving heart, such as Dante gives us, can be applied to the instinctive transports of an obedient soul before the authority that he judges to suit him: a quick judgment grants him to conceive that it will be good for him to serve that force conceived as useful and beneficial, whose order presages protection, justice, or victory. There is a taste in it of the beginning of a mysterious good. By what sign is it known? That is the great difficulty. Certain military leaders make themselves obeyed by genius, others by bravery, others by a type of mystical faith. The exterior and brilliant gifts of a Condé can add the magic of the example to it, but some generals carried on litters have radiated the same prestige.

Henry Fouqier, who was part of the \emph{Thousand} [who conquered Sicily as part of the unification of Italy], loved to tell that the aging Garibaldi inflamed his band by saying to them in a whisper, from the base of his carriage where rheumatism confined him, a simple: ``Gentlemen, move ahead!" So many passions of hope and confidence sleep in the human soul! Little suffices a little in order to make them rise from it, but this nothing is indispensable and no convention, no arrangement, no artifice of will can take the place of the first natural gift.

Authority is of the same order as virtue or genius or beauty.

The most learned cogs have never replaced authority \emph{born} …

The French of the tenth century had settled down around the race which, for a hundred years or more, had always defended them effectively. Where did that race come from, from what heaven had it fallen on the country? Immigrant Saxons or native rural lords or even descendants of bourgeois Parisians, scholarship does not stop discussing it. They do not discuss the authority acquired little by little by their fortuitous power nor the benefit of their dynasty nor its constant happiness.

It expresses for centuries a power or protection and increase, it represents everything that the heart and spirit of men, isolated or reunited, wait for, hope for, and believe about a true authority.

True authority is naturally wise; an insane authority is not conceivable. The idea of authority does not mean, indeed, only the power and the great power exercised by a man or a group of men, but all the more, it confines the knowledge of the object on which this power is exercised and applied. The more the authority believes, the more this knowledge itself is developed. The more authority is perfect, the more it assumes the clarity and the exactitude of that knowledge, and the more it proportions itself in it.

Authority would not be an eternal political necessity if, at the same time as that guiding instinct which constitutes the bottom of the soul of the leaders, there did not exist in the soul of the subjects and the citizens an instinct of obedience, ``spirit of following," said Richelieu, who is the living expression of the greatest interest of the crowds: to be governed and well governed, in a good sense, with firmness.

\paragraph{The Conditions of True Authority: the Education of the Leaders}

The development of what is called modern civilization tends increasingly favour material forces over moral forces.

If we rely on it in order to achieve social justice or the improvement of morals, we set ourselves up for considerable deceptions! That civilization equalizes neither wealth nor conditions: on the contrary, its complexity never ceases increasing the differences between men. It does not liberate: the authority of science and industry will tend instead to establish new races of slaves. Finally, far from pacifying and reconciling, its needs are so demanding that they seem to intersect at a right angle, to destroy or deny everything that is human.

Neither the game of supply and demand which continues capitalism, nor the principle of nationalities which created our armed peace, nor class warfare, by which the insurgent masses respond to the starvation brought by capitalism, would know how to spread in the modern world an ambiance of a sheepfold. We would be rather be pushed back among the antisocial wolves and pressured to live, by categories of class or race, according to the custom of the wolves. The inherited veneer of customs peels off little by little, the vestiges of general traditions disappear, and the statistics of criminality show what inevitably results from it.

Consider the progress of athletics (which could make, in a well ruled society, an admirable school of discipline and elegance), the passion (otherwise excellent in itself) of winning and prevailing in violent games, the new instruments created by science and those who apply them: this proliferation of ancient forces and these new means put in the service of powers without restraint only have to give their lessons: the moralist, with a thousand signs, sees brutality reborn.

As to the language of our contemporaries, I speak of the best, of those who are elevated, if not highly elevated, men and women, that it comes back to primitive onomatopoeia, if we ``let come to grips with life", forces, destinies.

In the regime of brutality, there are neither leaders nor orders, nor even order, without its necessary hierarchy which is lacking. The massacres of September had leaders. It is not orders that the executioner lacked from the Duke of Enghien or Bishop Darboy. It is not necessary to complain about our times in this regard.

Class differences are more marked than they have been for a half century, the arrogance and despotism of the authorities would rather be on the way to grandeur. What is lacking in the directing spirits is that light which is the sign of the right to lead. The leaders remain and their power increases, but those are barbaric leaders left to the impulses of passion or interest. They command, they lead, because their troops will it, but they command poorly and lead wrongly, the mistake being learned.

They themselves also are, still more than these proletarian masses for which they feign a very vivid interest, the really underprivileged.

The intellectual and moral treasure, whose task of collecting that inheritance belongs to them, was disdained and finally lost. Thus the spirit of liberal democracy used it, which disorganized the country from above; borrowing the language of progress, pretending to possess the promises of the future, it abandoned the sole instrument of progress, which is tradition, and the only seed of the future, which is the past.

The history of the Third Republic can suffice to show the harm that there can be to surrender legislation, armies, the economy, diplomacy, and all the forms of authority and influence to those spirits without direction and without culture, to those hearts without self-mastery and without dignity.

The symmetric history of conservative England, where everyone who governed and served in the top jobs had undergone the hard and long intellectual and moral preparation in the old universities, with a lot of Greek verses and Latin discourses, confirms how certain it is that the real happiness of the people depends on the good training of their leaders. The sword of the conqueror, the baton of the founder, even the pencil of the stockbroker, all these modalities of strength and cunning can and must bring about great benefits to the condition of having passed the necessary time under the strict rule of the educator. Everything that they take away from the rule, is not taken away from the rule nor from the authority which holds it: that is entrenched in the entire mass of the people; it is the nation and the human race who are the first dispossessed.

The lessening of the common intellectual and moral assets is a loss for everyone: the least will lose as much as the great.

They will even lose much more than the great, because that which perfects, refines, elevates the great, constitutes, to the profit of the others, the most precious, and often the only, guarantee against the abuses of power which greatness is precisely exposed to. Certain nuances of virtue and honor, certain beautiful persuasive accents of the voice which commands are the direct fruits of the sole education.

It is like that of religion.

Whoever said that a religion was necessary for the people spoke complete nonsense. Religion is necessary, education is necessary, a set of powerful brakes is necessary for the leaders of the people, for is advisors, its chiefs, due to the very function of direction and restraint that they are called to keep from it: if the passions of the human beast are to be feared for all, it will be to dread in the proportion that the beast will play with stronger powers and will be able to devastate a more extended field of action.

All liberty is not suitable in every State; each State depends on its historic antecedents and it geographic position like each man on his ancestors and his country. Salutary and tutelary dependencies, since they gave life, sustain it, and conserve it, and whoever rejects it, dies. Liberty varies with time and place, but this is no State which can last without a sovereign authority.

Therefore, if we want to speak with exactitude, it is not liberty which is general and necessary, an ecumenical, basic, and human right; it is authority.

\paragraph{The Exercise of Authority}

The property of power resembles other properties; it results from work, from work done and ``done well". Totally naked force can be applied for good and for evil, for construction, for destruction. When it did the good, when it constructed, it has merit from it, it has prestige and glory from it, it also sees the birth from it of this product which is called authority.

A power vacuum resembles the vacancy of a field. Whoever wants it, takes it; whoever is able to, holds onto it.

When there is a power vacuum it is, as Joan of Arc said, a great pity for the kingdom. And it is a great misfortune. To seize power in that case, if one has the strength, is simply an act of charity and humanity. A people needs a leader like a man needs bread. Not only, in such a hypothetical situation, is the right of first occupancy established, but there is a rigorous duty, a strict obligation for those who \emph{can} occupy. When citizens are threatened by the enemy, it is necessary to command them if one can do it. When disorder is in the street, it is necessary to restore it back to order if one has the means.

Power is not an idea, it is a fact, and one believes in this fact when it makes itself felt; all criticism by the world can do nothing against the strength of a conqueror.

Most moralizers, who have confused minds, judged that power corrupted the heart of man. When power is elevated and lasts, when it lasts a little, the effect is totally the opposite; the learning of responsibilities matures and their experience perfects, rather than ruins, them.



\flrightit{Posted on 2012-10-21 by Cologero }

\begin{center}* * *\end{center}

\begin{footnotesize}\begin{sffamily}



\texttt{Jason-Adam on 2012-11-02 at 15:02 said: }

Here is an interview deBenoist gave where he discusses Maurras in depth.

A case could be made from an Evolian perspective that Maurras did not go back far enough i.e. he defended l'ancien regime while forgetting that the actions of the French monarchy against the empire and papacy are what began the slide to destruction. Evola did say that the French monarchy brought on their own doom by their absolutist, anti-feudal and anti-Empire policies. 

Dante was also opposed to the entire Capetian house and saw them as nearly usurpers. (see Purgatorio)

Does anyone know if Evola actually wrote about Maurras ?


\hfill

\texttt{Cologero on 2012-11-02 at 22:22 said: }

I don't have time myself to find that passage. Where did Evola write that ``the actions of the French monarchy against the empire and papacy are what began the slide to destruction"? Then we will be in a position to discuss it.

When specifically did those actions occur? I will be making a case soon that the ``ancien regime" defined Feudal society through its understanding of the tri-functional division of society as pointed out by Dumezil. If the monarchy rebelled against the spiritual authority at some point then, yes, that would be indicative of the ``degeneration of castes". However, I just translated a text where Maurras explicitly said that the temporal leaders were to undergo a religious education (i.e., the spiritual authority).

Evola did say that he believed what well-bred men believed to be healthy and normal prior to the French Revolution, i.e., the ``ancien regime". That is the motivation behind the Maurras translations, viz., to highlight those healthy and normal beliefs.

Given this, I think Benoist's case falls apart.

I have just one complete article by Evola on Maurras from La Vita Italiana: ``The defeat and future of France according to Action Francaise". I'll try to get it translated before the end of the year.

In other books, there are some references in passing . I'll gather them together at some point in the near future.


\hfill

\texttt{Jason-Adam on 2012-11-03 at 01:47 said: }

1. I foolishly forgot to post the link to the interview, here it is

\url{http://www.alaindebenoist.com/pdf/entretien\_charles\_maurras.pdf}

2. Evola's anti French remarks are all throughout part II of Revolt against the Modern World. He specifically attacks Phillipe IV's attack of Pope Boniface VIII, also Francois I warring with Emperor Charles to the point of even allying with the anti-Christian Turks, Cardinal Richelieu's centralising policy which destroyed the independence of the local aristocracy, and then Louis XIV's continuing wars against Austria and the empire.

Also in Men among the Ruins Evola said Italy should move away from France and closer to Germany.

These are not my opinions, I am merely trying to provide Evola's thoughts as I understand them.


\hfill

\texttt{Cologero on 2012-11-04 at 11:23 said: }

Yes, Jason-Adam, that is the interview that ``inspired" the series on Maurras. Unfortunately, there are too many elements for a comment box: the diverse views of Maurras, Evola, and Benoist himself, not to mention the view from Tradition.

As you point out, discrete historical events over the course of centuries can be pointed out, as Evola did. Or else, we can look at these events from the point of view of Tradition: which were informed by a sense of order, and which by an opposition to order. To make things manageable, we therefore reject out of hand the notion of an ``Evolian" critique; we are only concerned about the critique of Tradition.

Specifically, then, a fundamental point is the degeneration of castes: how do Evola and Maurras see that. Of course, Evola was explicit about that, and he did not see the ``ancien regime" as disordered; quite to the contrary, as we have pointed out numerous times, Evola himself claimed his views were in harmony with those of well-bred men of that ancien regime. I would say the same for Maurras.

Benoist, on the other hand, rejects outright that regime, placing himself in opposition both to Evola and Maurras.

You do point out an important difference between Evola and Maurras, viz., their relationship to German thought. The latter was appalled by it. Maurras actually proposed a sort of greater Romanity, among all the nations speaking Romance languages. Hence, he hoped for a closer cooperation between France and Italy, leaving Germany out of it. Maurras saw a direct line from Athens to Rome to Paris. Evola, on the other hand, was obsessed with Germany and its philosophers.

So what next?

P.S. To those who don't read French: that interview is by a living writer under copyright protection; otherwise I would make it available in English for you.


\hfill

\texttt{Jason-Adam on 2012-11-04 at 22:13 said: }

Our discussion has reached an important highmark – Evola and Maurras are agreed upon which system of rule is best. We can also agree that restoring l'ancien regime, if it could be done, would result in restoring the Traditional form of western society. The major difference between Evola and Maurass is indeed the nature of Germany.

The question I have now is : on the matter of Germany and German philosophy, who is correct – Evola or Maurras ?


\hfill

\texttt{Cologero on 2012-11-04 at 23:52 said: }

Coincidentally, I have a post planned on German philosophy for mid-week. Ironically, Maurras ended up liking the Germans a bit too much, and that cost him his reputation and his freedom after WW II.

There is an important difference, also, in fundamental orientation. Evola and Maurras were both pagans enamoured of ancient Greece and Rome–although Maurras was more influenced by its classical aspects–and both had high regard for the Medieval period. Evola accepted Tradition as a metaphysical, supra-historical, given. Maurras was a Positivist and historian; his view of tradition was simply as an historical reality and his understanding of it came indirectly from his knowledge of Catholic France.


\hfill

\texttt{Logres on 2012-11-04 at 15:39 said: }

\url{http://distributistreview.com/mag/2012/08/lord-acton-tends-to-corrupt/}

Lord Acton's dictum has lead to a great deal of mischief in the Anglo ghetto.


\hfill

\texttt{Cologero on 2012-11-05 at 13:19 said: }

Judging from personal correspondence, there seems to be some confusion about this section. They seem to think that this is a justification for ``fascism". I responded that to me, it seems descriptive, not prescriptive. Actually, it can be read as the program for revolution as we wrote about recently. Revolutionary forces deliberately provoke public disorder precisely in order to create a power vacuum which they will attempt to fill. Since the public prefers order to chaos, they are usually welcomed … at first.

In the light of the previous section, remember that Maurras insisted that the temporal powers needed to be under the influence of religion. This is not true of modernity, following the American, French, and Russian revolutions which were overtly anti-religion to various degrees. That people who consider themselves pious or spiritual would favor modernity over the traditional thinkers who acknowledge the priority of the spiritual over the temporal, absolutely baffles me. They believe that modernity is the only ``loving" and sane system.

This should also be read in the context of the ``Social Contract" theories. Obviously, Maurras cannot accept that. The social contract theorists postulate that free individuals gather together to work out the mechanisms of power in their social group. Clearly that is not consistent with Maurras' understanding, since he claims that some group will assert power rather than wait for a social contract. The contract theorists presume the individual is supreme, whereas the traditional view is that society is prior to the individual.


\hfill

\texttt{Dominion on 2012-11-05 at 13:56 said: }

It should be noted, though, that not all religions are equal in this regard. Many people who subscribe to some religious principle or order have felt themselves justified to do great damage to their societies and others, to destroy a Traditional order and heritage (Ansar Dine in the great city of Timbuktu), and so on. If the religion is not grounded in Truth, then it can be as destructive as any ``atheist Stalinist."

The social contract can also be read in terms of different groups in society coming together, which is far closer to the way things really work (force isn't discounted). Preferably, the `contract' would then allow each member of each group the best way to self-development in their respective contexts and according to their abilities.

\hfill

\end{sffamily}\end{footnotesize}
